\scanpage{148}
Ceci implique que si $u$ est un isomorphisme, alors
$A(u)$ et $I(u)$ sont des isomorphismes inverses l'un
de l'autre. Enfin
\[
A(\sigma_g)=I(\sigma_g)=\text{identité dans } A(G)\quad(\text{si }g\in G). \tag*{(27)}
\]

Ces propriétés suffisent pour mettre sur la somme
directe $A$ des $A^n$ :
\[
\begin{aligned}
A^n&=A(\unsure{\mathfrak S|\mathcal S}_n)\quad(n\geq1)&\big|&& A^*&=\coprod_{n=0}^{\infty}A_{\unsure{n}}\\
A^0&=\mathbb Z
\end{aligned}\tag*{(28)}
\]
une structure \deleted{\unsure{de}} anneau gradué commutatif.
\inserted{(et de même sur les $A(G)^n=A^G(\mathfrak S_n)$, $A(G)^*=\coprod_{0}^{\infty}A(G)^n$}

D'ailleurs, on a, pour $\alpha\in A^p$, $\beta\in A^q$
\[
A(u_{pq})(\alpha*\beta)=\frac{(p+q)!}{p!\,q!}\alpha\square\beta
\quad\in A^{\mathfrak S_p\times\mathfrak S_q} \tag*{(29)}
\]
d'où, si on veut une relation dans $A^{\unsure{*|\bullet}}$ lui-même
\[
A(u_{p,p})(\alpha*\beta)=\frac{(\deleted{p+q}\inserted{2p})!}{p!\,q!}\alpha\beta \tag*{(30)}
\]
qui relie la multiplication dans les $A^n$ \deleted{\unsure{avec}}
à la multiplication ``globale'' $*$.

De plus, la formule 16 montre qu'on
a un hom canonique
\[
\left\{\begin{aligned}
&T\underset{\unsure{\cdot}}{=}T_G:A(G)\longrightarrow1+\widehat{A(G)^*}{}^+ \deleted{\unsure{z_+}}\\
&T(\alpha)=\sum T^n(\alpha)\deleted{\unsure{t^n}}\qquad(T^0(\alpha)=1)
\end{aligned}\right.\tag*{(31)}
\]
défini par
\[
T(\gamma(M))=\sum T^n(M)\deleted{\unsure{t^n}} \tag*{(32)}
\]
On a donc
\[
T^n(\alpha+\beta)=\sum_{p+q=n}T^p(\alpha)*T^q(\beta) \tag*{(33)}
\]
On déduit de (21)
\[
T^n(\alpha\beta)=T^n(\alpha)T^n(\beta) \tag*{(33)}
\]
(qui, compte tenu de (29), exprime ``presque''
% archive sheet number: 147
% printed form (vertical paper watermark, reverse face)
\unsure{GENTINE} \tikz[baseline=-.2ex,scale=.17]{\draw (0,0)--(1.6,0)--(1.8,.33)--(.5,.33)--(.32,.13)--(.92,.13);} TYPEWRITER \tikz[baseline=-.4ex,scale=.15]{\foreach \a in {0,45,...,315}{\draw (\a:.52) circle (.22);} \draw (0,0) circle (.21);} \unsure{ARGENT}
